% Gerado por regressao/06_tabelas_latex.R (projeto Modelagem). Não editar à mão.
\begin{landscape}
\begin{table}[htbp]
\caption{Robustez R15 — sem as variáveis do Censo (jovens, idosos, urbanização)}\label{tab:robustez_R15}
\centering
\footnotesize
\setlength{\tabcolsep}{3pt}
\begin{tabular}{@{}>{\raggedright\arraybackslash}p{\dimexpr\linewidth-594pt\relax}>{\centering\arraybackslash}p{60pt}>{\centering\arraybackslash}p{60pt}>{\centering\arraybackslash}p{60pt}>{\centering\arraybackslash}p{60pt}>{\centering\arraybackslash}p{60pt}>{\centering\arraybackslash}p{60pt}>{\centering\arraybackslash}p{60pt}>{\centering\arraybackslash}p{60pt}>{\centering\arraybackslash}p{60pt}@{}}
\hline
\textbf{Variável} & \textbf{\hspace{0pt}Despesa total} & \textbf{\hspace{0pt}Saúde e Saneamento} & \textbf{\hspace{0pt}Educação e Cultura} & \textbf{\hspace{0pt}Habitação e Urbanismo} & \textbf{\hspace{0pt}Assistência Social} & \textbf{\hspace{0pt}Transporte} & \textbf{\hspace{0pt}Administração} & \textbf{\hspace{0pt}Agricultura} & \textbf{\hspace{0pt}Comunicações} \\ \hline
\multicolumn{10}{@{}l}{\textit{Modelo A (ciclo eleitoral)}} \\
Ano eleitoral & 277,669\textsuperscript{***} & 53,095\textsuperscript{***} & 50,014\textsuperscript{*} & 126,479\textsuperscript{***} & 14,082\textsuperscript{***} & 40,661\textsuperscript{***} & 6,814 & 0,110 & $-$0,793\textsuperscript{***} \\
 & (101,989) & (10,712) & (30,316) & (31,498) & (4,302) & (10,036) & (10,768) & (2,119) & (0,230) \\
Ano pré-eleitoral & 161,903\textsuperscript{*} & 0,814 & 82,310\textsuperscript{**} & 33,086\textsuperscript{**} & 9,543\textsuperscript{**} & 5,341 & 10,318 & 3,687 & 0,665\textsuperscript{***} \\
 & (93,496) & (11,921) & (39,071) & (16,263) & (4,649) & (7,197) & (8,209) & (3,464) & (0,150) \\
\hline
Observações & 70.713 & 70.570 & 70.594 & 69.613 & 70.529 & 54.449 & 70.597 & 64.498 & 12.360 \\
Municípios & 5.568 & 5.568 & 5.568 & 5.565 & 5.568 & 5.207 & 5.568 & 5.437 & 1.896 \\
R\textsuperscript{2} & 0,836 & 0,752 & 0,694 & 0,298 & 0,372 & 0,171 & 0,275 & 0,159 & 0,030 \\[6pt]
\multicolumn{10}{@{}l}{\textit{Modelo B (coalizões)}} \\
Prefeito na coalizão & 31,000\textsuperscript{***} & 4,393\textsuperscript{*} & 0,224 & 4,096 & 0,687 & 1,332 & 7,319\textsuperscript{***} & $-$0,954 & $-$0,125 \\
do governador & (6,300) & (2,355) & (3,166) & (2,799) & (0,720) & (2,694) & (2,728) & (1,054) & (0,301) \\
Prefeito na coalizão & $-$10,974 & 7,683\textsuperscript{**} & 2,033 & 0,064 & 0,979 & $-$15,707\textsuperscript{***} & 0,488 & $-$5,766\textsuperscript{***} & 0,467 \\
do presidente & (7,617) & (3,121) & (3,973) & (4,133) & (0,997) & (3,234) & (3,818) & (1,404) & (0,410) \\
\hline
Observações & 70.713 & 70.570 & 70.594 & 69.613 & 70.529 & 54.449 & 70.597 & 64.498 & 12.360 \\
Municípios & 5.568 & 5.568 & 5.568 & 5.565 & 5.568 & 5.207 & 5.568 & 5.437 & 1.896 \\
R\textsuperscript{2} & 0,551 & 0,349 & 0,283 & 0,106 & 0,143 & 0,071 & 0,136 & 0,056 & 0,022 \\
\hline
\end{tabular}
\fonte{Elaboração própria.}
\nota{Modelo A: efeito fixo de município, erros-padrão de Driscoll-Kraay. Modelo B: efeitos fixos de município e de ano, erros-padrão agrupados por município. Erros-padrão entre parênteses. \textsuperscript{***}p$<$0,01; \textsuperscript{**}p$<$0,05; \textsuperscript{*}p$<$0,1.}
\end{table}
\end{landscape}
